\documentclass{article}

\usepackage[main, final]{neurips_2026}

\usepackage[utf8]{inputenc}
\usepackage[T1]{fontenc}
\usepackage{hyperref}
\usepackage{url}
\usepackage{booktabs}
\usepackage{amsfonts}
\usepackage{amsmath}
\usepackage{nicefrac}
\usepackage{microtype}
\usepackage{xcolor}
\usepackage{graphicx}
\usepackage{multirow}

\title{Not All Calibrations Are Equal:\\Language Bias in LLM Pruning}

\author{%
  Minguhn Kim \\
  M.S.\ Candidate, BEVIL\\
  GSST, POSTECH\\
  \texttt{minguhn@postech.ac.kr} \\
  \And
  Jihun Park \\
  M.S.\ Candidate, ISDS\\
  GSCST, POSTECH\\
  \texttt{jpark42@postech.ac.kr} \\
}

\begin{document}
\maketitle

\begin{abstract}
Wanda prunes large language models by scoring each weight as the product of its magnitude and the $\ell_2$-norm of the corresponding input activation, computed from a small calibration set.
The original paper demonstrates robustness to the \emph{number} of calibration samples but does not test sensitivity to the domain or language of the calibration data.
We provide evidence that the activation norm acts as a statistical fingerprint of the calibration distribution, and that this fingerprint determines which model capabilities survive pruning once sparsity is high enough.
On LLaMA-3-8B with four calibration sources (English, Korean, code, random tokens) at 50--70\% unstructured sparsity, we find: (1)~at 50\%, calibration source spreads WikiText-2 perplexity by 6.27 points but all Wanda variants cluster within 1.2 percentage points on KoBEST HellaSwag (Korean), consistent with the original robustness finding; (2)~at 70\%, English-calibrated Wanda loses 12 percentage points more Korean capability than Korean-calibrated Wanda, while both damage English knowledge equally; (3)~WikiText-2 perplexity, the original paper's primary metric, consistently favors English calibration and cannot detect the Korean-capability damage.
Within statistical uncertainty, Korean calibration dominates on downstream task accuracy: it preserves Korean at no cost to English.
\end{abstract}

%------------------------------------------------------------------
\section{Introduction}
%------------------------------------------------------------------

Deploying large language models (LLMs) efficiently requires compression.
Among compression techniques, unstructured pruning zeroes out individual weights to reduce effective model size.
Wanda~\citep{sun2024wanda}---Pruning by \textbf{W}eights \textbf{and} \textbf{A}ctivations---is a widely adopted LLM pruning baseline owing to its simplicity: for each weight $w_{ij}$, the importance score is
\begin{equation}
  S_{ij} = |w_{ij}| \cdot \|x_j\|_2,
  \label{eq:wanda}
\end{equation}
where $x_j$ is the $j$-th input feature activation collected over a small calibration set.
Weights with the lowest scores are zeroed out.
No gradient computation or weight update is required, and the method runs in minutes even for 65B-parameter models.

The original evaluation demonstrates robustness to the \emph{number} of C4~\citep{raffel2020t5} samples (English web text), varying from 128 up to 2048 with little change in WikiText-2 perplexity.
But it does not test sensitivity to the \emph{domain} or \emph{language} of the calibration data.
Since $\|x_j\|_2$ is computed from the calibration set, it encodes which features are active in that data.
Features that are dormant during calibration---such as features specific to Korean when calibrating on English---receive low activation norms, low importance scores, and are pruned first.
We call this the \emph{calibration fingerprint}: the activation statistics encode the calibration distribution and bias which weights are pruned.

Several recent works support this concern: calibration data choice significantly affects pruning quality~\citep{ji2025calibration}, even moderate sparsity substantially harms multilingual performance~\citep{choenni2025multilingual}, and reconstruction-based pruning can overfit its calibration data~\citep{shin2024calibration}.
Yet no prior work directly tests whether the calibration fingerprint produces \emph{asymmetric} capability loss across languages, or whether the standard evaluation metric (WikiText-2 perplexity) can detect such loss.

We test this on LLaMA-3-8B~\citep{grattafiori2024llama3} with four calibration sources at 50--70\% unstructured sparsity.
First, at 50\% sparsity, calibration source clearly separates WikiText-2 perplexity but does \emph{not} separate KoBEST HellaSwag (Korean) or MMLU (English)---the original paper's robustness finding holds at this operating point.
Second, at 70\% sparsity, English-calibrated Wanda loses 12 percentage points (pp) more Korean capability than Korean-calibrated Wanda, while both damage English knowledge equally (within 0.24\,pp on MMLU at 50\%, where both scores are well above chance).
The asymmetry is single-sided: Korean calibration preserves Korean without sacrificing English.
On downstream task accuracy, it is the dominant calibration choice within statistical uncertainty.
Third, WikiText-2 perplexity at 70\% favors English calibration (108.37 vs.\ 138.24), the \emph{opposite} direction from the Korean-capability result---the original paper's primary metric actively recommends the calibration choice that maximizes minority-language damage.

Following the Track~2 stress-testing protocol, we also document reproducibility issues in the original codebase, including a hardcoded calibration buffer that silently prevents the paper's own sensitivity sweep from running correctly.

%------------------------------------------------------------------
\section{Background and Related Work}
\label{sec:background}
%------------------------------------------------------------------

\textbf{Wanda} prunes each linear layer independently by computing $S_{ij}$ (Eq.~\ref{eq:wanda}) for all weights and zeroing those below a target sparsity ratio~\citep{sun2024wanda}.
The key insight over magnitude-only pruning is that a weight connected to a high-activation feature should be treated as more important, even if its magnitude is small---an idea grounded in the observation that LLMs develop outlier features with disproportionately large activations~\citep{dettmers2022llmint8}.
The calibration set is used solely for a single forward pass to collect activation statistics; no weight updates are performed.

\textbf{Magnitude pruning} scores weights by $|w_{ij}|$ alone, requiring no calibration data.
\textbf{SparseGPT}~\citep{frantar2023sparsegpt} solves a layer-wise reconstruction problem that updates remaining weights after pruning, making it more expensive but generally stronger.
Both serve as baselines: magnitude as a calibration-free lower bound, SparseGPT as a calibration-dependent upper bound.

\textbf{Calibration sensitivity in pruning.}
Ji et al.~\citep{ji2025calibration} showed that calibration data is crucial to post-training pruning quality, especially at high sparsity, and that data more similar to the pre-training distribution yields better results---but tested only English evaluation.
Choenni \& Titov~\citep{choenni2025multilingual} proposed M-Wanda and showed that even moderate sparsity ratios substantially harm multilingual performance, using mixed multilingual calibration rather than isolating single-language calibration effects.
Shin et al.~\citep{shin2024calibration} demonstrated that reconstruction-based pruning can overfit its calibration data, and proposed self-generated calibration to mitigate this.
Our contribution is the direct isolation of the calibration-language variable: we compare English-only vs.\ Korean-only calibration and evaluate on both English and Korean benchmarks.
The result is a single-sided asymmetry invisible to perplexity-based evaluation.

%------------------------------------------------------------------
\section{Experimental Setup}
\label{sec:setup}
%------------------------------------------------------------------

\textbf{Model.}\quad LLaMA-3-8B~\citep{grattafiori2024llama3}, chosen because it scores above random chance on KoBEST HellaSwag (Korean), a prerequisite for measuring Korean-capability loss after pruning.

\textbf{Calibration sources.}\quad We use 128 samples per source, matching the original Wanda default:
\emph{C4}~\citep{raffel2020t5} (English web text, the original paper's default),
\emph{MC4-ko} (Korean subset of multilingual C4),
\emph{The Stack} (source code, a non-natural-language domain),
and \emph{random tokens} (sampled uniformly from the vocabulary; a structureless control).

\textbf{Evaluation.}\quad All metrics are computed with lm-evaluation-harness~\citep{gao2024lmeval}.
\emph{WikiText-2 perplexity} ($\downarrow$): the original Wanda paper's primary metric, measuring English language-modeling fluency.
\emph{MMLU}~\citep{hendrycks2021mmlu} ($\uparrow$, 5-shot accuracy): English world knowledge and reasoning (14,042 items; stderr $\approx$\,0.4\,pp).
\emph{KoBEST HellaSwag}~\citep{jang2022kobest} ($\uparrow$, 0-shot accuracy, 4-way MCQ): Korean commonsense reasoning (stderr $\approx$\,2.2\,pp; chance level 25\%)---the capability we expect to be most sensitive to English-only calibration.
GSM8K (math reasoning) was planned but excluded entirely due to CUDA kernel failures during chain-of-thought generation that affected all model variants (see \S\ref{sec:repro}).

\textbf{Sparsity levels.}\quad 50\%, 60\%, and 70\% unstructured sparsity, applied per output neuron as in the original paper.

\textbf{Compute.}\quad 8$\times$ NVIDIA RTX 6000 Ada (48\,GB each), CUDA 12.6, PyTorch 2.6.0, Transformers 5.5.3.
All experiments total 14 distinct pruning runs on LLaMA-3-8B.

%------------------------------------------------------------------
\subsection{Reproducibility Issues}
\label{sec:repro}
%------------------------------------------------------------------

Before running stress tests, we attempted to reproduce original results using the public codebase.\footnote{\url{https://github.com/locuslab/wanda}}
Three issues are especially consequential:

\emph{(1)~Hardcoded calibration buffer.}
The function \texttt{prepare\_calibration\_input} allocates the calibration tensor with a hardcoded \texttt{torch.zeros((128, ...))}, ignoring the \texttt{-{}-nsamples} argument.
With $\texttt{nsamples}>128$ the first-layer hook overflows with an \texttt{IndexError}; with $\texttt{nsamples}<128$ the buffer contains trailing zeros read as valid calibration.
This silently prevents the paper's own sensitivity sweep (their Table~17) from running as described.

\emph{(2)~Transformers kwargs mismatch.}
The calibration catcher reads \texttt{kwargs['position\_ids']} directly, but Transformers 4.40+ may pass \texttt{position\_ids=None} or introduce \texttt{position\_embeddings} as a separate kwarg, causing a \texttt{KeyError}.
Our fix captures all kwargs into a dict and re-emits them at the layer call site.

\emph{(3)~CUDA kernel failures during evaluation.}
Two distinct failure modes affected our evaluation.
First, GSM8K's chain-of-thought generation (1,319 requests) crashed with \texttt{CUDA error: unspecified launch failure} for all model variants including dense, making GSM8K unrecoverable.
Second, MMLU's loglikelihood evaluation (56,168 requests) crashed intermittently for \emph{Wanda-pruned} LLaMA-3-8B specifically, but not for magnitude- or SparseGPT-pruned models on the same hardware; the error originates in the SwiGLU MLP, suggesting numerical instability from the Wanda per-output mask.
Setting \texttt{-{}-eval\_batch\_size\,1} resolves the MMLU crash for most configurations.
MMLU cells at 60\% sparsity for Wanda were not measured; the thesis comparison is captured at 50\% and 70\%.

Three additional issues (torch version blocking, C4 dataset config rename, OPT architecture mismatch) are documented in our fork.\footnote{Available upon request; will be released publicly after the course evaluation period.}

%------------------------------------------------------------------
\section{Results}
\label{sec:results}
%------------------------------------------------------------------

\subsection{Calibration Source at 50\% Sparsity}

Table~\ref{tab:calibration} holds model and sparsity fixed and varies the calibration source for Wanda, comparing against magnitude pruning (calibration-free) and SparseGPT (calibration-dependent).

\begin{table}[t]
  \caption{Effect of calibration source on Wanda at 50\% unstructured sparsity (LLaMA-3-8B). WikiText-2 perplexity separates calibration sources clearly; KoBEST HellaSwag does not.}
  \label{tab:calibration}
  \centering
  \small
  \begin{tabular}{@{}lccc@{}}
    \toprule
    Method (calibration) & WT-2 ppl\,$\downarrow$ & MMLU\,$\uparrow$ & KoBEST HS\,$\uparrow$ \\
    \midrule
    Dense (no pruning) & 5.54 & 65.40 & 56.60 \\
    \midrule
    Magnitude (no calibration) & 310.71 & 26.79 & 36.80 \\
    SparseGPT (C4) & 8.52 & 53.73 & 52.20 \\
    \midrule
    Wanda (C4, English) & 9.06 & 48.59 & 51.60 \\
    Wanda (MC4-ko, Korean) & 10.09 & 48.83 & 51.00 \\
    Wanda (The Stack, code) & 10.18 & 46.31 & 52.20 \\
    Wanda (random tokens) & 15.33 & 30.11 & 51.00 \\
    \bottomrule
  \end{tabular}
\end{table}

WikiText-2 perplexity follows a clear ordering: C4 (9.06) $<$ MC4-ko (10.09) $\approx$ The Stack (10.18) $<$ random (15.33), a 6.27-point spread consistent with the activation norm encoding the calibration distribution.
On MMLU, the three natural-language calibrations are within 2.5\,pp of each other (C4 and MC4-ko within 0.24\,pp), while random tokens collapse to 30.11\%.
On KoBEST HellaSwag, all four Wanda variants cluster in a 1.2\,pp window (51.00--52.20\%), well within the $\pm$2.2\,pp stderr.
At this sparsity, calibration source does not measurably affect Korean capability.

Magnitude pruning collapses on LLaMA-3-8B (perplexity 310.71, MMLU at chance), unlike the moderate degradation reported for LLaMA-2 in the original paper.
This prevents our pre-registered test of whether mismatched-calibration Wanda could fall below magnitude pruning: even Wanda with random tokens beats this baseline by 14.2\,pp on KoBEST.

\subsection{Sparsity $\times$ Calibration Interaction}

Table~\ref{tab:sparsity} sweeps sparsity for the two calibration sources of primary interest (C4 and MC4-ko).
Figure~\ref{fig:kobest} plots KoBEST HellaSwag across sparsity for all methods.

\begin{table}[t]
  \caption{Wanda sparsity sweep: English (C4) vs.\ Korean (MC4-ko) calibration on LLaMA-3-8B. Dashes indicate MMLU values not measured at 60\% due to CUDA evaluation failures (\S\ref{sec:repro}). The KoBEST gap between C4 and MC4-ko grows from $-$0.60\,pp at 50\% to +12.00\,pp at 70\%.}
  \label{tab:sparsity}
  \centering
  \small
  \begin{tabular}{@{}llccc@{}}
    \toprule
    Calibration & Sparsity & WT-2 ppl\,$\downarrow$ & MMLU\,$\uparrow$ & KoBEST HS\,$\uparrow$ \\
    \midrule
    \multirow{3}{*}{C4 (English)}
      & 50\% & 9.06 & 48.59 & 51.60 \\
      & 60\% & 23.07 & --- & 43.00 \\
      & 70\% & 108.37 & 26.36 & 29.00 \\
    \midrule
    \multirow{3}{*}{MC4-ko (Korean)}
      & 50\% & 10.09 & 48.83 & 51.00 \\
      & 60\% & 32.62 & --- & 47.00 \\
      & 70\% & 138.24 & 26.21 & 41.00 \\
    \bottomrule
  \end{tabular}
\end{table}

\begin{figure}[t]
  \centering
  \includegraphics[width=0.85\linewidth]{kobest_sparsity.pdf}
  \caption{KoBEST HellaSwag accuracy across sparsity levels. Wanda (C4) and Wanda (MC4-ko) start together at 50\% and diverge at 70\%, with a 12\,pp gap. SparseGPT (C4) degrades more gracefully than Wanda (C4). Magnitude pruning collapses below random chance by 60\%.}
  \label{fig:kobest}
\end{figure}

The KoBEST gap between C4 and MC4-ko calibration grows monotonically with sparsity: $-$0.60\,pp at 50\% (noise), +4.00\,pp at 60\% (emerging), +12.00\,pp at 70\% (well outside statistical noise: the difference stderr is $\sqrt{2}\times2.2\approx3.1$\,pp, giving ${\approx}$\,3.9$\sigma$).
Equivalently, the KoBEST drop from 50\% to 70\% is $-$22.60\,pp for C4 but only $-$10.00\,pp for MC4-ko---English calibration suffers more than twice the Korean-capability loss under the same sparsity increase.

MMLU tells the opposite story.
At 50\%, C4 and MC4-ko produce 48.59\% and 48.83\% respectively (difference 0.24\,pp, well within the $\pm$0.4\,pp stderr at $n$\,=\,14k).
At 70\%, both collapse to near chance: 26.36\% and 26.21\% (difference 0.15\,pp).
At both sparsities where MMLU was measured, the two calibrations damage English knowledge equally.

%------------------------------------------------------------------
\section{Analysis}
\label{sec:analysis}
%------------------------------------------------------------------

\subsection{The Calibration Fingerprint}

The WikiText-2 perplexity ordering at 50\% (Table~\ref{tab:calibration}) is consistent with $\|x_j\|_2$ carrying a distributional signature: in-domain English calibration (C4) produces the best English perplexity; out-of-domain calibration (MC4-ko, The Stack) adds ${\sim}$1\,ppl; structureless calibration (random tokens) adds 6.27\,ppl.
One anomaly stands out: random-token calibration collapses MMLU to 30.11\% (near chance) while leaving KoBEST at 51.00\% (indistinguishable from other calibrations).
This likely reflects KoBEST's limited statistical power ($\pm$2.2\,pp stderr) rather than a true difference in vulnerability; moderate damage would be undetectable at this sample size.

\subsection{Sparsity-Dependent Emergence}

At 50\% sparsity, the calibration fingerprint is measurable in perplexity but does not affect KoBEST: all four Wanda variants produce statistically indistinguishable Korean scores (Table~\ref{tab:calibration}).
On MMLU, the two natural-language calibrations (C4 and MC4-ko) are equally indistinguishable, though code and random calibrations degrade English knowledge.
The pruning budget is loose enough that features survive regardless of their calibration-derived importance.
At 70\%, the budget tightens and the fingerprint becomes a triage mechanism---the KoBEST gap grows from noise at 50\% to a clear signal at 70\%, with the effect already emerging at 60\% (Table~\ref{tab:sparsity}).
The calibration-source effect on Korean capability is therefore a sparsity-dependent phenomenon: invisible at the operating point tested in the original paper, decisive at aggressive sparsity.

\subsection{Single-Sided Asymmetry}

The asymmetry between English and Korean capability loss is the central finding.
Korean calibration (MC4-ko) preserves Korean capability without sacrificing English: at 50\%, where both MMLU scores are well above chance (48.83 vs.\ 48.59), the difference is 0.24\,pp---conclusively null.
English calibration (C4) destroys Korean capability at high sparsity yet gains nothing on English over Korean calibration: at 70\%, MMLU is 26.36 vs.\ 26.21 (0.15\,pp difference) while KoBEST is 29.00 vs.\ 41.00 (12.00\,pp difference).

We attribute the single-sidedness to an asymmetry in the calibration data itself.
Korean web text naturally contains English tokens---borrowed words, mixed script, Latin-alphabet brand names---so MC4-ko calibration activates both Korean- and English-relevant features.
English web text contains essentially no Korean, so C4 calibration activates only English-relevant features and assigns near-zero importance to Korean-specific ones.
We did not directly verify this by inspecting per-feature activation norms, but the explanation is consistent with all observed data.

On downstream task accuracy, MC4-ko matches or exceeds C4 within statistical uncertainty at every sparsity tested: it matches C4 on MMLU (differences of 0.24\,pp and 0.15\,pp, both well inside stderr) and exceeds it on KoBEST at 60\% and 70\%.
This makes MC4-ko the dominant calibration choice for any deployment that values Korean capability---it costs nothing on English.
Note that this dominance holds on downstream accuracy, not on WikiText-2 perplexity, where C4 is consistently better (9.06 vs.\ 10.09 at 50\%; 108.37 vs.\ 138.24 at 70\%).

\subsection{Evaluation Blindness}

The most practically important result is the perplexity--capability dissociation at 70\% sparsity.
WikiText-2 perplexity says C4 calibration is 29.87 points better than MC4-ko (108.37 vs.\ 138.24).
KoBEST says MC4-ko calibration is 12.00\,pp better than C4 (41.00 vs.\ 29.00).
The two metrics give opposite recommendations.

A practitioner following the standard evaluation pipeline---calibrate on English, evaluate on perplexity---would choose C4 at every sparsity and never discover the Korean-capability damage.
This is not a failure of perplexity as a metric \emph{per se}; WikiText-2 perplexity faithfully measures English fluency.
The failure is in using an English-only metric to evaluate the effect of calibration-data choice, when the risk is precisely that calibration biases pruning toward the language of the calibration data.
The evaluation is structurally blind to the capability damage it should be measuring.

\subsection{Secondary Findings}

\textbf{Magnitude collapse on LLaMA-3-8B.}\quad
Magnitude pruning at 50\% produces WikiText-2 perplexity of 310.71 on LLaMA-3-8B, compared to ${\sim}$15 on LLaMA-2-7B reported in the original paper.
On LLaMA-3-8B, calibration-free pruning collapses catastrophically---similar to the OPT/BLOOM behavior documented by~\citet{sun2024wanda}.
The original paper's favorable comparison of Wanda against magnitude pruning does not hold for LLaMA-3-8B.

\textbf{SparseGPT graceful degradation.}\quad
SparseGPT degrades more gracefully than Wanda across sparsity on this model (perplexity 37.52 vs.\ 108.37 at 70\%; KoBEST 32.40 vs.\ 29.00).
However, Wanda with Korean calibration (MC4-ko) at 70\% outperforms SparseGPT with English calibration (C4) on KoBEST by 8.6\,pp (41.00 vs.\ 32.40), despite being ${\sim}$300$\times$ faster.
This comparison confounds algorithm and calibration choices, but the practical lesson is clear: choosing the right calibration for a simple method can matter more than choosing a sophisticated algorithm with default calibration.

%------------------------------------------------------------------
\section{Limitations}
\label{sec:limitations}
%------------------------------------------------------------------

Each experiment uses a single random seed; variance across seeds is not reported.
The effect sizes on KoBEST at 70\% (12\,pp, ${\approx}$\,3.9$\sigma$ on the difference stderr) are large relative to statistical noise, but the exact numbers may shift with different seeds.
All experiments use LLaMA-3-8B; we do not test generalization to other model families (Qwen, Mistral) or model sizes.
Korean is the only non-English language tested; the mechanism (calibration fingerprint biasing pruning) suggests this should generalize, but we have not verified it.
GSM8K evaluation was lost entirely to CUDA failures, and MMLU has gaps at 60\% sparsity for Wanda variants, leaving the per-task fragility analysis incomplete.
KoBEST HellaSwag has a small test set ($\pm$2.2\,pp stderr); effects below ${\sim}$5\,pp---including the predicted calibration advantage at 50\% sparsity---cannot be reliably detected at this sample size.

%------------------------------------------------------------------
\section{Conclusion}
%------------------------------------------------------------------

We stress-tested Wanda's calibration robustness by varying calibration domain and language rather than quantity.
Our results are consistent with the activation norm acting as a calibration-distribution fingerprint.
At 50\% sparsity, this fingerprint creates a measurable perplexity ordering across calibration sources but does not affect downstream task accuracy.
At 60--70\%, it becomes a capability triage mechanism.
The resulting asymmetry is single-sided and grows with sparsity: Korean calibration preserves Korean capability (KoBEST gap grows from noise at 50\% to +12\,pp at 70\%) while English knowledge is unaffected (MMLU within 0.24\,pp at every sparsity).
WikiText-2 perplexity, a standard pruning evaluation metric, cannot detect this damage and actively favors the calibration choice that maximizes it.
For any deployment valuing minority-language capability, in-language calibration is the dominant strategy.

\begin{thebibliography}{11}

\bibitem[Sun et~al.(2024)]{sun2024wanda}
Sun, M., Liu, Z., Bair, A., \& Kolter, J.\,Z.
\newblock A simple and effective pruning approach for large language models.
\newblock In \emph{ICLR}, 2024.

\bibitem[Frantar \& Alistarh(2023)]{frantar2023sparsegpt}
Frantar, E. \& Alistarh, D.
\newblock {SparseGPT}: Massive language models can be accurately pruned in one shot.
\newblock In \emph{ICML}, 2023.

\bibitem[Raffel et~al.(2020)]{raffel2020t5}
Raffel, C., Shazeer, N., Roberts, A., Lee, K., Narang, S., Matena, M., Zhou, Y., Li, W., \& Liu, P.\,J.
\newblock Exploring the limits of transfer learning with a unified text-to-text transformer.
\newblock \emph{JMLR}, 21(140):1--67, 2020.

\bibitem[Grattafiori et~al.(2024)]{grattafiori2024llama3}
Grattafiori, A. et~al.
\newblock The {L}lama 3 herd of models.
\newblock \emph{arXiv preprint arXiv:2407.21783}, 2024.

\bibitem[Dettmers et~al.(2022)]{dettmers2022llmint8}
Dettmers, T., Lewis, M., Belkada, Y., \& Zettlemoyer, L.
\newblock {LLM.int8()}: 8-bit matrix multiplication for transformers at scale.
\newblock In \emph{NeurIPS}, 2022.

\bibitem[Ji et~al.(2025)]{ji2025calibration}
Ji, Y., Xiang, Y., Li, J., Xia, Q., Li, P., Duan, X., Wang, Z., \& Zhang, M.
\newblock Beware of calibration data for pruning large language models.
\newblock In \emph{ICLR}, 2025.

\bibitem[Choenni \& Titov(2025)]{choenni2025multilingual}
Choenni, R. \& Titov, I.
\newblock {M-Wanda}: Improving one-shot pruning for multilingual {LLM}s.
\newblock In \emph{EMNLP}, 2025.

\bibitem[Shin et~al.(2024)]{shin2024calibration}
Shin, S., Park, W., Lee, J., \& Lee, N.
\newblock Rethinking pruning large language models: Benefits and pitfalls of reconstruction error minimization.
\newblock In \emph{EMNLP}, 2024.

\bibitem[Gao et~al.(2024)]{gao2024lmeval}
Gao, L., Tow, J., Abbasi, B., Biderman, S., Black, S., DiPofi, A., Foster, C., Golding, L., Hsu, J., Le Noac'h, A., Li, H., McDonell, K., Muennighoff, N., Ociepa, C., Phang, J., Reynolds, L., Schoelkopf, H., Skowron, A., Sutawika, L., Tang, E., Thite, A., Wang, B., Wang, K., \& Zou, A.
\newblock A framework for few-shot language model evaluation.
\newblock \emph{Zenodo}, 2024.
\newblock \url{https://github.com/EleutherAI/lm-evaluation-harness}.

\bibitem[Hendrycks et~al.(2021)]{hendrycks2021mmlu}
Hendrycks, D., Burns, C., Basart, S., Zou, A., Mazeika, M., Song, D., \& Steinhardt, J.
\newblock Measuring massive multitask language understanding.
\newblock In \emph{ICLR}, 2021.

\bibitem[Jang et~al.(2022)]{jang2022kobest}
Jang, M., Kim, D., Kwon, D.\,S., \& Davis, E.
\newblock {KoBEST}: {K}orean balanced evaluation of significant tasks.
\newblock In \emph{COLING}, 2022.

\end{thebibliography}

\end{document}
